Practitioners often assume that greedy (temperature-0) decoding makes an LLM's output a deterministic function of its prompt. We test whether this holds when the same prompt is served in different batches. We run 1,700 byte-identical prompts from five tasks through \vllm under static batch sizes $B\in\{1,4,16,64\}$, a same-size reshuffle of batch-mates, and exact repeats. We use ten configurations: five bf16 models (0.5B--3B, \qwen and \llama) plus fp16, fp32 and AWQ variants. Moving from sequential to batched decoding changes the exact output of 20--34\% of prompts for bf16 models, and of 47--76\% of long chain-of-thought outputs. This is far above the 3--5\% we hypothesized. Batch size itself is not the driver. Reshuffling batch-mates at fixed $B=64$ causes nearly as much divergence, and repeating an identical schedule causes none in all ten configurations. Every bf16 divergence starts at a near-tie: about 45\% at exact top-2 ties and at least 91\% within one bf16 rounding step. Higher precision reduces divergence on the same weights from 33.5\% (bf16) to 7.3\% (fp16) and 2.8\% (fp32). Task type has no effect once output length is controlled, and model size shows no trend. Mean accuracy changes by at most 1.8\pp and never significantly after Holm correction, yet 3--5\% of items, and up to 16\% of chain-of-thought multiple-choice items, flip between correct and incorrect. Single-run bf16 evaluations therefore carry hidden item-level noise, and per-item reproducibility requires a fixed batch schedule, batch-invariant kernels, or higher precision.
